\footnotesize


\noindent Zu Band VIII, 1:
\setlength\LTleft{0pt}\setlength\LTright{0pt}
\begin{longtable}{p{36mm}p{92mm}}

S.~\href{\LSBviiiURL LSB_VIII_1.pdf#page=147}{102}, Z.~3\textendash5 & \textit{statt} Im oberen \lbrack...\rbrack\ tang rad \textit{lies} Im oberen Viertel folgende Notiz von fremder Hand, die inhaltlich mit dem Diagramm \lbrack\textit{Fig.~1}\rbrack\ zusammenhängt: \textit{72 + 5 (id} \textlangle\textit{est}\textrangle\ \textit{77) ad rd ita} \lbrack\textit{ut}\rbrack\ \textit{72 ad tang} \lbrack/\rbrack\ \textit{rad}\\%

\hangindent=3.3em 
S.~\href{\LSBviiiURL LSB_VIII_1.pdf#page=148}{103}, Erläuterung zu Z.~3\newline\qquad 
Thevenotiani & \textit{ergänze} 
Siehe zudem \textsc{Ders.}, \textit{Machine nouvelle pour la conduite des eaux, pour les bâtimens, pour la navigation, et pour la pluspart des autres arts}, Paris 1665.\\%

S.~\href{\LSBviiiURL LSB_VIII_1.pdf#page=179}{134}, nach Z.~8 & \textit{ergänze:} \textit{E\textsuperscript{3}}\enskip \textit{Notitia Opticae promotae}, in \textsc{B.~Spinoza}, \textit{Adnotationes ad Tractatum theologico politicum}, hrsg. von C.\,G.\;von~Murr, den Haag 1802, S.~24\textendash29.\\

S.~\href{\LSBviiiURL LSB_VIII_1.pdf#page=179}{134}, Z.~10 & \textit{statt} \textit{E\textsuperscript{3}} \textit{lies} \textit{Ü}\\%

S.~\href{\LSBviiiURL LSB_VIII_1.pdf#page=192}{147}, Z.~3 & \textit{statt} Leibn. Marg. 163 \textit{lies} Hannover GWLB Nm-A 163\\%

S.~\href{\LSBviiiURL LSB_VIII_1.pdf#page=675}{630}, \textit{Fig.~1} & \textit{Ergänze Erl.}: Gemäß den Angaben im Text muss das große Rad 18 Zähne besitzen.\\%

S.~\href{\LSBviiiURL LSB_VIII_1.pdf#page=720}{675}, Z.~3 & \textit{statt} Leibn. Marg. 163 \textit{lies} Nm-A 163\\%

\end{longtable}

\vspace{2ex}


\noindent Zu Band VIII, 2:
\setlength\LTleft{0pt}\setlength\LTright{0pt}
\begin{longtable}{p{36mm}p{92mm}}

S.~\href{\LSBviiiURL LSB_VIII_2.pdf#page=92}{44}, Z.~1\textendash3	&	%
\textit{streiche} \lbrack p.~256\rbrack\ \lbrack...\rbrack\ posse:\\

S.~\href{\LSBviiiURL LSB_VIII_2.pdf#page=92}{44} 		&			%
\textit{streiche das Diagramm Fig.~3 samt Erläuterung, da die fraglichen Zeichen nicht auf Leibniz zurückzuführen sind, sondern bereits Bestandteil der gedruckten Vorlage sind}\\

S.~\href{\LSBviiiURL LSB_VIII_2.pdf#page=93}{45}	&	%
\textit{statt} \lbrack\textit{Fig.~4}\rbrack\ \textit{lies} \lbrack\textit{Fig.~3}\rbrack\\

S.~\href{\LSBviiiURL LSB_VIII_2.pdf#page=434}{386}, Z.~13 & \textit{streiche} Erit ergo \lbrack...\rbrack\ \textit{ER} \textit{RQ}.\\

S.~\href{\LSBviiiURL LSB_VIII_2.pdf#page=437}{389} & \textit{ergänze Anmerkung zu Fig.~3}: Erit ergo vis in composita ratione brachiorum \textit{FE} \textit{GE} et rectarum \textit{ER} \textit{RQ}.\\

S.~\href{\LSBviiiURL LSB_VIII_2.pdf#page=714}{666}, Z.~2 & \textit{statt} Tasterin \textit{lies} Laßerin\\

S.~\href{\LSBviiiURL LSB_VIII_2.pdf#page=714}{666}, Z.~4 & \textit{statt} Tasterin \textit{lies} Laßerin\\

\end{longtable}

\vspace{2ex}


\noindent Zu Band VIII, 3:
\setlength\LTleft{0pt}\setlength\LTright{0pt}
\begin{longtable}{p{36mm}p{92mm}}

S.~\href{\LSBviiiURL LSB_VIII_3.pdf#page=12}{xii}, Z.~17	&	%
\textit{statt} Berichtigungen
\textit{lies} Berichtigungen und Ergänzungen\\

S.~\href{\LSBviiiURL LSB_VIII_3.pdf#page=15}{xv}, Z.~22	&	%
\textit{statt} ritter.bbaw.de
\textit{lies} leibniz-berlin.bbaw.de\\ 

S.~\href{\LSBviiiURL LSB_VIII_3.pdf#page=16}{xvi}, Z.~15
&	%
\textit{statt} erschienen Stücken
\textit{lies} erschienenen Stücke\\

S.~\href{\LSBviiiURL LSB_VIII_3.pdf#page=16}{xvi}, Z.~29
&	%
\textit{ergänze} Gedankt sei Olga Fedorova, die Dimitri Bayuk bei der Aufnahme von Stücken half.\\

S.~\href{\LSBviiiURL LSB_VIII_3.pdf#page=41}{7}, Z.~5		&				%
\textit{statt} Streichungen
\textit{lies} Anstreichungen \\

S.~\href{\LSBviiiURL LSB_VIII_3.pdf#page=70}{36}, Z.~16	&				%
\textit{statt} ispius
\textit{lies} ipsius \\

S.~\href{\LSBviiiURL LSB_VIII_3.pdf#page=77}{43}, Z.~18 	&				%
\textit{statt} naturalem\textit{TAQ}
\textit{lies} naturalem~\textit{TAQ} \\

S.~\href{\LSBviiiURL LSB_VIII_3.pdf#page=87}{53}, Z.~10 	&				%
\textit{statt} vibr.\textit{MLP}
\textit{lies} vibr.~\textit{MLP} \\

\hangindent=3.5em 
S.~\href{\LSBviiiURL LSB_VIII_3.pdf#page=257}{223}, Variante \textit{E\textsuperscript{1}} zu\newline\qquad Z.~4\textendash S.~224, Z.~1 
& 	%
\textit{statt} aliqasque
\textit{lies} aliasque \\

S.~\href{\LSBviiiURL LSB_VIII_3.pdf#page=262}{228}, Z.~1 	&				%
\textit{statt} quantum
\textit{lies} tum \\

S.~\href{\LSBviiiURL LSB_VIII_3.pdf#page=264}{230}, Z.~7 	&				%
\textit{statt} \textit{G}trabem
\textit{lies} \textit{G}~trabem \\

S.~\href{\LSBviiiURL LSB_VIII_3.pdf#page=266}{232}, zu \textit{Fig. 5a}	&		%
\textit{statt} \textit{Bl.~1~r\textsuperscript{o}}
\textit{lies} \textit{Bl.~2~r\textsuperscript{o}} \\

\hangindent=3.5em 
S.~\href{\LSBviiiURL LSB_VIII_3.pdf#page=293}{259}, Erläuterung zu\newline\qquad Z.~11\,f.	
&	%
\textit{statt} SieheN.
\textit{lies} Siehe N. \\

S.~\href{\LSBviiiURL LSB_VIII_3.pdf#page=404}{370}, Anmerkung 	&			%
\textit{statt} versu\lbrack s\rbrack\ 3\textsuperscript{\lbrack a\rbrack} \lbrack...\rbrack\ ictus\textsuperscript{\lbrack b\rbrack} 
\textit{lies}  versu\lbrack s\rbrack\textsuperscript{\lbrack a\rbrack} 3\textsuperscript{\lbrack b\rbrack} \lbrack...\rbrack\ ictus\textsuperscript{\lbrack c\rbrack}
\textit{und im dazugehörigen Marginalienapparat statt} \textsuperscript{\lbrack a\rbrack}~\glqq3\grqq\ bezeichnet \lbrack...\rbrack\ Körpers. \textsuperscript{\lbrack b\rbrack}~versu \textit{L ändert Hrsg.}
\textit{lies} \textsuperscript{\lbrack a\rbrack}~versu \textit{L ändert Hrsg.} \textsuperscript{\lbrack b\rbrack}~\glqq3\grqq\ bezeichnet \lbrack...\rbrack\ Körpers.
\\

S.~\href{\LSBviiiURL LSB_VIII_3.pdf#page=405}{371}, Z.~3 	&			%
\textit{statt} Seiten
\textit{lies} Seite \\

S.~\href{\LSBviiiURL LSB_VIII_3.pdf#page=433}{399}, Z.~21\,f. 	&	%
\textit{statt} seit Huygens' Abreise aus Paris im Juli 1676
\textit{lies} seit Huygens' krankheitsbedingter Abreise aus Paris im Juli 1676 (siehe dazu seinen Brief an Leibniz von Ende Juni 1676: \textit{LSB} III,~1 N.~86, S.~428\,f.) \\

S.~\href{\LSBviiiURL LSB_VIII_3.pdf#page=434}{400}, Z.~7 	&	%
\textit{statt} N.~40,
\textit{lies} N.~40. \\

S.~\href{\LSBviiiURL LSB_VIII_3.pdf#page=483}{449}, Z.~1	& 	%
\textit{statt} Gründen
\textit{lies} Gründe \\

S.~\href{\LSBviiiURL LSB_VIII_3.pdf#page=539}{505}, Z.~4 	&	%
\textit{statt} \lbrack149~r\textsuperscript{o}\rbrack\
\textit{lies} Bl.~149~r\textsuperscript{o} \\ 

S.~\href{\LSBviiiURL LSB_VIII_3.pdf#page=633}{599}, Z.~6 	&	%
\textit{statt} bei
\textit{lies} an\\

S.~\href{\LSBviiiURL LSB_VIII_3.pdf#page=649}{615}, 
\href{\LSBviiiURL LSB_VIII_3.pdf#page=650}{616},
\href{\LSBviiiURL LSB_VIII_3.pdf#page=652}{618},
\href{\LSBviiiURL LSB_VIII_3.pdf#page=653}{619}
&	%
\textit{beachte die inkonsistente Zeilenzählung}\\

S.~\href{\LSBviiiURL LSB_VIII_3.pdf#page=675}{641}, Z.~1	&	%
\textit{statt} eadem,etiam
\textit{lies} eadem, etiam\\

S.~\href{\LSBviiiURL LSB_VIII_3.pdf#page=675}{641}, variante zu Z.~9	&	%
\textit{statt} distantiam\textit{L}
\textit{lies} distantiam~\textit{L}\\

S.~\href{\LSBviiiURL LSB_VIII_3.pdf#page=702}{668}, Z.~6	&	%
\textit{statt} \sout{\textit{lieffe}} \textit{lies} \textit{\sout{ließe}}\\

S.~\href{\LSBviiiURL LSB_VIII_3.pdf#page=702}{668}, Z.~29	&	%
\textit{statt} 1675\textendash1677
\textit{lies} 1695\textendash1697\\

S.~\href{\LSBviiiURL LSB_VIII_3.pdf#page=713}{679}, Z. 16  	& 	%
\textit{statt} (LH XXXV 11, 18B Bl.~207~v\textsuperscript{o}\textendash210~v\textsuperscript{o}) 
\textit{lies} (LH XXXV 11, 18C Bl.~207~v\textsuperscript{o}\textendash210~v\textsuperscript{o}) \\

S.~\href{\LSBviiiURL LSB_VIII_3.pdf#page=715}{681}, Z.~32\textendash33
&				%
\textit{streiche} in Bl.~98 \lbrack...\rbrack\ Das Zeichen\\

S.~\href{\LSBviiiURL LSB_VIII_3.pdf#page=715}{681}, Z.~33
&				%
\textit{statt} Bl.~96
\textit{lies} Bl.~94\\	

S.~\href{\LSBviiiURL LSB_VIII_3.pdf#page=739}{705}, Z.~6 	&	%
\textit{statt} Kustode
\textit{lies} Kustoden \\

S.~\href{\LSBviiiURL LSB_VIII_3.pdf#page=739}{705}, Z.~13 	&	%
\textit{statt} im den
\textit{lies} in dem \\

S.~\href{\LSBviiiURL LSB_VIII_3.pdf#page=789}{755}, Z.~29	&	%
\textit{statt} Vgl. a.a.O., S. 37.
\textit{lies} Vgl. das Diagramm in \textsc{I.~Newton}, \textit{Principia}, lib.~I, sect.~II, prop.~I (London 1687, S.~37).\\

S.~\href{\LSBviiiURL LSB_VIII_3.pdf#page=797}{763}, Anmerkung 	& 			%
\textit{statt} \textit{An Rand}
\textit{lies} \textit{Am Rand} \\

S.~\href{\LSBviiiURL LSB_VIII_3.pdf#page=854}{820}, Z.~1 	&				%
\textit{statt} les corps
\textit{lies} corps \\

S.~\href{\LSBviiiURL LSB_VIII_3.pdf#page=875}{841}, Z.~4 	&	%
\textit{statt} ein Absatz
\textit{lies} fünf Zeilen \\



\end{longtable}



\vfill					%
\noindent
Hinweise und Korrekturen werden dankbar entgegengenommen unter
\href{mailto:leibniz8@bbaw.de}{leibniz8@bbaw.de}
\vspace*{3\baselineskip}	%
